\begin{textsample}{POS Dim 3 – go – Score 30.00 – go/go\_inf\_25.txt}  \label{ex:f3_pos_050}
3.5.2.
Organização da Ação Pedagógica A \textbf{criança} para descobrir a si mesma e o mundo precisa, desde \textbf{bebê}, ter a possibilidade de frequentar diferentes espaços, ter contato com a natureza e realizar atividades ao ar livre.
Em média, \textbf{instituições} de tempo integral atendem as \textbf{crianças}, de 9 ( nove ) a 10 ( dez ) horas por \textbf{dia}, ficando boa parte do tempo em ambientes fechados.
As de tempo regular, com 4 ( quatro ) horas por \textbf{dia}, geralmente, destinam o horário do recreio e do parquinho como os momentos possíveis para as \textbf{crianças} ficarem fora da sala.
Essa realidade \textbf{demonstra} a necessidade de mudar concepções e repensar os usos dos espaços, de forma que a \textbf{criança} extrapole os muros da \textbf{instituição}, tendo como perspectiva, o entorno e a comunidade na qual está inserida, para que possa conhecer outros locais e vivenciar novas possibilidades de aprendizagens ( TIRIBA, 2010 ).
Tanto na cidade como no campo podem ser explorados espaços, como : praças, \textbf{parques}, clubes, quintais de vizinhos com jardins ou hortas, campos de \textbf{futebol}, quadras cobertas, ginásios, laboratórios de informática, cooperativas, indústrias, órgãos de saneamento básico, matas, reservas ambientais, morros, rios, descampados, cachoeiras, pastos, chácaras e fazendas para observar plantações, currais, chiqueiros, galinheiros, produção de queijo, de farinha, de rapadura, de polvilho etc.
Esses espaços de forma integrada ao planejamento do professor e da \textbf{instituição} educacional, enquanto coletivo, podem permitir a \textbf{criança} : - sentir o calor do sol, o vento, o orvalho da manhã, os pingos da chuva ; - observar e elaborar questões sobre animais, plantas, fenômenos naturais, transformação de materiais ; - explorar \textbf{areia}, grama, lama, terra, ar, água, gravetos, pedras ; - experimentar cheiros, \textbf{texturas}, \textbf{sons}, sabores, cores ; - desenvolver noção espacial ao fazer relações como perto ou longe ; - descrever os trajetos que serão realizados, informando se passa por ruas, avenidas, rios, becos, trieiros, mata-burros etc. ; - identificar os diferentes tipos de vegetação, edificações de moradias que podem ser visualizadas ao longo do caminho.
Dessa forma, a \textbf{percepção} e compreensão desses outros espaços como educativos, potencializa a articulação com a comunidade, a ideia de pertencimento a um determinado grupo social e permite o conhecimento das histórias locais e dos sujeitos mais marcantes da comunidade, bairro, vilarejo, cidade.
Na \textbf{instituição} educacional, mesmo em espaços \textbf{pequenos}, ainda é possível realizar estudos junto à natureza, ou pensar com a \textbf{criança} em alternativas para a construção de jardins ou hortas suspensas, minhocários, composteiras, terráreos, sementeiras, definição de locais adequados para armazenamento de lixos que podem ser reciclados ou reaproveitados, a fim de compreenderem determinados processos da natureza, assim como, desenvolver atitudes mais conscientes na relação com o ambiente.
Espaços mais estruturados como pesque-pague, museus, zoológicos, planetário, \textbf{bibliotecas}, universidades e faculdades constituem também lugares que possibilitam pesquisas, observações, levantamento de hipóteses, debates e questionamentos.
A ida a esses diferentes locais deve ser caracterizada como momento de vivência e de estudo.
Para tanto, tem que ser devidamente planejada junto às \textbf{crianças} para definirem horário, \textbf{dia}, organização da turma ou da \textbf{instituição}, os materiais que serão levados, o que será observado, os \textbf{registros} que serão feitos ( escrita, vídeo, fotografia, desenho ) e por quem.
É importante, também, reafirmar os \textbf{cuidados} que todos devem ter com o espaço, por exemplo : sempre colocar o lixo no lugar apropriado ; preservar os seres vivos encontrados ; coletar material para pesquisa apenas quando houver necessidade.
Posteriormente, na \textbf{instituição}, é importante retomar o que foi observado e debatido nas visitas, avaliando enquanto grupo a ação realizada, os conhecimentos aprendidos e as questões que podem ser elaboradas para dar continuidade na ampliação do tema/assunto que está sendo abordado em sequências didáticas, projetos de trabalho, temáticos e/ou investigativos.
Outra forma de possibilitar a diversificação de ações e a ampliação de conhecimentos às \textbf{crianças} é a \textbf{instituição} educacional estabelecer parcerias com programas de saúde, unidades itinerantes de oftalmologia, odontologia, dermatologia, combate ao Aedes aegypti.
A promoção de eventos, mostras científicas ou \textbf{feira} de ciências pela \textbf{instituição} educacional, também, possibilita a socialização com as famílias e com a comunidade, dos conhecimentos e saberes apropriados pelas \textbf{crianças} a partir de projetos e de experimentos.
Os materiais, assim como o espaço, são fundamentais na promoção de aprendizagens e desenvolvimento, necessitando de investimentos financeiros tanto para a sua aquisição, quanto para a formação de professores no sentido de utilizá -los de forma contextualizada com as \textbf{crianças}, explorando ao máximo suas possibilidades, no desenvolvimento do pensamento científico e na elaboração de generalizações.
Os materiais permitem à \textbf{criança} : a observação de animais e de plantas, de fenômenos da natureza, de relações de causa e efeito, de mudanças nas paisagens ; a exploração de elementos da natureza e de suas transformações ; a manipulação de objetos e a \textbf{percepção} de noções espaciais, temporais e de quantidade ; a realização de experimentos ; dentre tantas outras possibilidades, a fim de \textbf{apoiar} a \textbf{criança} na construção de suas próprias teorias e na compreensão do mundo físico e do mundo sociocultural.
Abaixo, são elencados alguns objetos e materiais que podem favorecer a efetivação dos objetivos de aprendizagens e desenvolvimento deste campo de experiências, tais como : - lanternas, lupas e microscópios para observação e exploração de espaços, insetos, folhas etc. ; - ampulhetas, relógios digitais e analógicos para \textbf{medição} de tempo ; - instrumentos de medidas variados, como fitas métricas, trenas, metros, réguas, copos medidores, xícaras, \textbf{colheres}, balanças de precisão etc. - atlas, mapas e globos terrestres para conhecimento das diferentes representações espaciais ; - moedas e cédulas de \textbf{brinquedo} para compreensão do sistema monetário ; - jogos de mesa, como dominó, memória, ludo etc. ; - objetos de tamanhos, \textbf{texturas}, cores, formas diversas, e ; - celulares, máquinas \textbf{fotográficas} para \textbf{registros}.
Esses são só alguns exemplos, vários outros, próprios do \textbf{dia} a \textbf{dia}, como argila, água, anilina, terra, cones, retalhos de \textbf{tecido}, prendedores, utensílios \textbf{domésticos} etc., devem ser ainda pensados pelas \textbf{instituições}, de acordo com a realidade e com o contexto de cada comunidade educacional.
No que se referem às relações sociais e espaço temporais as experiências da \textbf{criança} podem ser ampliadas e diversificadas com : - exploração e manipulação de diferentes objetos e materiais etc. ; - realização de labirintos e túneis para as \textbf{crianças} experimentarem diferentes movimentos e posições - correr, saltar, pular, rolar, andar, agachar, deitar etc. ; - reconhecimento de diferentes lugares da sua moradia e da \textbf{instituição} que frequenta a partir do uso cotidiano ; - localização de variados elementos -móveis, \textbf{brinquedos}, materiais - a partir de um ponto de referência em situações de \textbf{brincadeiras} e jogos ; - conhecimento de várias formas de representar o espaço, as convencionais, como, os mapas e o globo terrestre e as não convencionais como os desenhos ; - uso de vocabulário adequado pelo professor na indicação de localizações - à direita, à esquerda, em cima, embaixo etc. - e na identificação de formas planas e tridimensionais ; - utilização de ampulhetas em situações de jogos que necessitam de marcação de tempo e \textbf{calendário} para \textbf{dias}, \textbf{meses} e ano ; - construção de linha do tempo da vida das \textbf{crianças} utilizando fotografias ou desenhos ; - exploração de relações entre objetos e espaços por meio dos jogos de construção, com blocos de materiais, tipos e tamanhos variados e outros \textbf{brinquedos} como, bonecas(os), \textbf{carrinhos}, entre outros.
Quanto ao planejamento do tempo, é importante destacar que os professores precisam prever tempo suficiente para a \textbf{criança} em diversas situações, pensar, elaborar suas respostas, observar, narrar suas conclusões, individualmente, em \textbf{pequenos} ou grandes grupos, de acordo com o ritmo pessoal de cada uma.
Essa ação favorece conhecer a forma de pensar de outras \textbf{crianças}, comparando estratégias e raciocínio utilizados.
Ter tempo para falar, expressar ideias, opiniões, hipóteses e \textbf{curiosidades} é fundamental para a \textbf{criança} desenvolver a fala, a \textbf{autoconfiança}, o poder de argumentação e o pensamento científico.
É necessário destacar que a regularidade, a continuidade e a apresentação de um mesmo conhecimento em contextos diferentes, com recursos e estratégias metodológicas variadas, possibilitam uma melhor compreensão pela \textbf{criança}.
Não é possível entender algo se ele for abordado uma única vez ou de forma esporádica.
Ter sequência e cadência nas ações pedagógicas realizadas nas \textbf{instituições} de Educação \textbf{Infantil} é fundamental para aprendizagens mais significativas. 4.
As Transições na Educação \textbf{Infantil}

% matched lemmas: apoiar, areia, autoconfiança, bebê, biblioteca, brincadeira, brinquedo, calendário, carro, colher, criança, cuidado, curiosidade, demonstrar, dia, doméstico, feira, fotográfico, futebol, infantil, instituição, medição, mês, parque, pequeno, percepção, registro, som, tecido, textura
\end{textsample}
